% \chapter{ANNEXURE}
\chapter{Engineering Considerations, Challenges and Remedies}

\section{Design Constraints}
\begin{itemize}
    \item \textbf{Ethical Considerations (WK9):} The framework handles sensitive personal information (e.g., student IDs, legal names, and academic records). To comply with ethical practice, data minimization principles are strictly applied: no plain-text credentials or transcripts are persistently written onto the public ledger. Full data consent protocols are implied, and encryption passphrases are kept securely client-side to prevent non-consensual tracking or data privacy violations.
    \item \textbf{Environmental Considerations (WK5):} High computation models and public blockchain proof-of-work (PoW) loops typically present deep carbon-footprint concerns. To mitigate this energy demand, this system adopts a private/permissioned strategy and utilizes a local test network configuration alongside a lightweight containerized stack (Node.js/Python), which minimizes continuous idle mining loops and excessive carbon output.
    \item \textbf{Social Constraints (WK7):} Centralized academic verification processes frequently pose geographical, administrative, and financial barriers for students seeking remote or global employment background screening. The system addresses this by providing a lightweight, responsive Web interface that requires no specific cryptographic wallet background knowledge from end-users, promoting open and equitable academic verification access across diverse student demographics.
    \item \textbf{Sustainability on Environmental and Societal Context:} Structurally, the smart contract relies on an append-only state machine where only the unique file hash (32 bytes) is anchored. This ensures data scalability and low-gas resource footprints on the cloud network while providing a lifelong, permanent verification standard that guarantees credential validity even if the original issuing academic center ceases operations.
\end{itemize}

\section{Relevant Challenges Addressed}
\begin{itemize}
    \item \textbf{The Pre-Validation Integrity Dilemma:} Standard document registries store files sequentially without inspecting contents beforehand, which creates a critical risk where an unverified or initially fraudulent certificate is permanently validated on-chain. The system remedies this by implementing a precise Python-driven ``pre-validation'' layer that screens text structures and layout footprints through Explainable AI before any block anchoring transaction occurs.
    \item \textbf{Algorithmic Transparency vs. Model Opaque Blocks:} Standard deep learning checkers operate as a black box, outputting a generic binary score that lacks legal or professional accountability. This research overcomes this by deploying Perplexity and Burstiness metrics that map out exactly which paragraphs exhibit robotic properties, transforming a blind classification into a human-interpretable report.
\end{itemize}

\section{Addressed Knowledge Profile Attributes}
\begin{itemize}
    \item \textbf{WK2 (Data Analysis \& formal CS aspects):} Applied conceptually based numerical statistics, text tokenization, cross-entropy modeling, and cosine vector similarities to execute advanced document indexing.
    \item \textbf{WK5 (Environmental Impacts \& Sustainability):} Implemented off-chain document hashing paired with localized consensus testing to minimize broad-scale infrastructure data inflation and excessive power costs.
    \item \textbf{WK9 (Professional Ethics \& Responsibilities):} Managed cryptographic permissions (Issuer, Verifier, Admin role hierarchies) inside Solidity code to prevent identities from being counterfeited or unethically altered.
\end{itemize}

\section{4. Addressed Complex Engineering Problems (CEP) \& Activities (CEA)}
\begin{itemize}
    \item \textbf{WP1 \& WP3 (Depth of Knowledge \& Abstract Thinking Required):} The integration requires combining non-linear linguistic models (Python/Transformers) with deterministic state ledgers (Solidity), demanding advanced abstract modeling to coordinate asynchronous processes without introducing single points of configuration failure.
    \item \textbf{EA1 \& EA2 (Range of Resources \& Optimal Interactions):} Effectively unified independent, heterogeneous technologies—serverless cloud databases (Neon PostgreSQL), localized blockchain compilers (Hardhat), vector distance indexes (pgvector), and NLP analysis pipelines—into a clean execution dashboard.
\end{itemize}

\newpage

